\chapter{Complexity — Formulas}
%%% generated by the blueprint pipeline from TCSlib.Complexity.CircuitComplexity.Formulas
\section{Overview}

Declarations of the Lean module \texttt{TCSlib.Complexity.CircuitComplexity.Formulas}.

\section{Declarations}

\begin{definition}[List of literals]
\lean{LitList}
\statementsource{Vad12}{pdf-bf2ce5422088-p032-b002}
\statementsource{arora-barak-ch13-circuit-lowerbounds}{pdf-52929dbaa962-p002-b002}
\statementsource{Analysis-of-Boolean-Functions-by-Ryan-ODonnell}{
  pdf-ffbe9c7b3306-p006-b002,
  pdf-ffbe9c7b3306-p093-b004
}
\label{LitList}
\leanok
\uses{Literal}
For a natural number $n$, a literal list is a finite list of literals on $n$ variables.
\end{definition}

\begin{definition}[Clause]
\lean{Clause}
\statementsource{arora-barak-ch13-circuit-lowerbounds}{pdf-52929dbaa962-p002-b002}
\statementsource{Analysis-of-Boolean-Functions-by-Ryan-ODonnell}{pdf-ffbe9c7b3306-p094-b004}
\label{Clause}
\leanok
\uses{LitList}
For a natural number $n$, a clause is a finite list of literals on $n$ variables.
\end{definition}

\begin{definition}[Depth-two formula representation]
\lean{Depth2}
\statementsource{Vad12}{pdf-bf2ce5422088-p032-b002}
\statementsource{arora-barak-ch13-circuit-lowerbounds}{
  pdf-52929dbaa962-p001-b005,
  pdf-52929dbaa962-p002-b002
}
\statementsource{Analysis-of-Boolean-Functions-by-Ryan-ODonnell}{
  pdf-ffbe9c7b3306-p006-b002,
  pdf-ffbe9c7b3306-p093-b004,
  pdf-ffbe9c7b3306-p094-b004
}
\label{Depth2}
\leanok
\uses{LitList}
For a natural number $n$, a depth-two formula representation is a finite list of terms,
each of which is a finite list of literals on $n$ variables.
\end{definition}

\begin{definition}[Width of a depth-two formula]
\lean{Depth2.width}
\statementsource{arora-barak-ch13-circuit-lowerbounds}{pdf-52929dbaa962-p002-b002}
\statementsource{Analysis-of-Boolean-Functions-by-Ryan-ODonnell}{
  pdf-ffbe9c7b3306-p014-b012,
  pdf-ffbe9c7b3306-p094-b002,
  pdf-ffbe9c7b3306-p094-b004
}
\label{Depth2.width}
\leanok
\uses{Depth2, LitList.width}
For a natural number $n$, the width of a depth-two formula representation is the largest
number of literals in any of its terms, with width $0$ when the representation has no
terms.
\end{definition}

\begin{theorem}[Evaluation of the empty term]
\lean{Term.eval_nil}
\label{Term.eval_nil}
\leanok
\uses{Term.eval}
\difficulty{1}
For any natural number $n$ and any Boolean assignment $x : \text{Fin}(n) \to
\{\mathrm{false},\mathrm{true}\}$, the empty term evaluates to true.
\end{theorem}

\begin{theorem}[Evaluation of a term with a leading literal]
\lean{Term.eval_cons}
\label{Term.eval_cons}
\leanok
\uses{Literal, Literal.eval, Term, Term.eval}
\difficulty{1}
Let $n$ be a natural number, let $l$ be a literal on $n$ variables, let $t$ be a list of
literals on those variables, and let $x$ be a Boolean assignment to the variables. The
term formed by placing $l$ at the beginning of $t$ evaluates at $x$ to the conjunction
of the value of $l$ at $x$ and the evaluation of $t$ at $x$.
\end{theorem}

\begin{theorem}[Evaluation of the empty clause]
\lean{Clause.eval_nil}
\label{Clause.eval_nil}
\leanok
\uses{Clause.eval}
\difficulty{1}
For any natural number $n$ and any assignment $x : \mathrm{Fin}\,n \to
\{\mathrm{false},\mathrm{true}\}$, the empty clause evaluates to false.
\end{theorem}

\begin{theorem}[Evaluation of a Clause with a First Literal]
\lean{Clause.eval_cons}
\label{Clause.eval_cons}
\leanok
\uses{Clause, Clause.eval, Literal, Literal.eval}
\difficulty{1}
Let $n$ be a natural number, let $l$ be a positive or negated literal on one of $n$
Boolean variables, let $c$ be a finite list of such literals, and let $x$ assign a
Boolean value to each variable. The clause formed by placing $l$ before $c$ evaluates to
the Boolean disjunction of the value of $l$ and the evaluation of $c$ on $x$.
\end{theorem}

\begin{theorem}[Evaluation of a Constructed DNF]
\lean{DNF.eval_mk}
\label{DNF.eval_mk}
\leanok
\uses{DNF.eval, Term, Term.eval}
\difficulty{1}
Let $n$ be a nonnegative integer, let $ts$ be a finite list of terms on $n$ Boolean
variables, and let $x$ be a Boolean assignment to those variables. The DNF whose terms
are $ts$ evaluates to true at $x$ exactly when at least one term in $ts$ evaluates to
true at $x$.
\end{theorem}

\begin{theorem}[Evaluation of a constructed CNF]
\lean{CNF.eval_mk}
\label{CNF.eval_mk}
\leanok
\uses{CNF.eval, Clause, Clause.eval}
\difficulty{1}
Let $n$ be a natural number, let $c_1,\ldots,c_r$ be a list of clauses on $n$ variables,
and let $x\in\{0,1\}^n$. The CNF formula formed from this list evaluates to true on $x$
exactly when every clause in the list evaluates to true on $x$.
\end{theorem}

\begin{theorem}[Width of a Constructed DNF]
\lean{DNF.width_mk}
\label{DNF.width_mk}
\leanok
\uses{DNF.width, Depth2.width, Term}
\difficulty{1}
Let $n$ be a nonnegative integer and let $ts$ be a finite list of terms on $n$ Boolean
variables. The width of the DNF with terms $ts$ equals the maximum number of literals in
any term of $ts$, with the maximum taken to be $0$ when $ts$ is empty.
\end{theorem}

\begin{theorem}[Width of a constructed CNF]
\lean{CNF.width_mk}
\label{CNF.width_mk}
\leanok
\uses{CNF.width, Clause, Depth2.width}
\difficulty{1}
Let $n$ be a natural number and let $c_1,\ldots,c_r$ be a list of clauses on $n$
variables. The width of the CNF formula formed from this list equals the maximum clause
width in the list, with width $0$ when the list is empty.
\end{theorem}
